The throughline of this report is a single, stubborn empirical fact:
on the reference 44M model, recurrence helps enormously within its
trained horizon and not at all outside it. Everything else here is an
attempt to understand, exploit, or work around that fact.

The architecture chapters show \emph{why} the model looks the way it
does: sparse routing concentrates learning signal per expert and
demands careful load balancing at small batch sizes; grouped-query
attention is a serving-time optimization whose training-time footprint
is negligible at this scale; and the looped structure converts
parameters into effective depth at the cost of gradient-conditioning
constraints that favour small, near-isometric recurrent blocks.

The training chapters propose how to do better next time: curricula
that order TinyStories by measurable complexity, and distillation
protocols that supervise the recurrent dynamics themselves rather than
only their final outputs. The evaluation chapters map the territory
around the model: which test-time scaling axes are viable below a
billion parameters, how quantization interacts with sparsity, what the
loops might mechanistically be computing, and how the field should
benchmark recurrent small models at all.

The most important next step is the one the centerpiece demands:
training with randomized loop counts to test whether extrapolation can
be \emph{taught} rather than hoped for---no experiment in this report
tests that yet. Until then, the honest summary of this programme is
that test-time loop scaling, on the evidence of one 44M run, offers no
free gains beyond the trained depth---and establishing that is itself
progress.
